% ============================================================
% FUENTES: Notas Biomate 1; raw_plain_nonlinear_dynamics.txt.
% Strogatz (2015), cap. 8; Murray (2002); Edelstein-Keshet (2007).
% ============================================================
\section{Bifurcación de Hopf}

En los capítulos de la Parte III estudiamos las bifurcaciones en flujos unidimensionales (saddle-node, transcrítica y pitchfork). En todas ellas, la pérdida de estabilidad de un equilibrio ocurría necesariamente a través de un autovalor real que cruzaba por cero ($\lambda = 0$).

Al expandir nuestro análisis a dos dimensiones, emerge un mecanismo de inestabilidad cualitativamente superior: un par de autovalores complejos conjugados puede cruzar el eje imaginario con velocidad no nula. Este fenómeno, descubierto por Aleksandr Andronov en la década de 1930 y popularizado por Eberhard Hopf en 1942, se conoce como la \textbf{bifurcación de Hopf} (o de Andronov--Hopf)~\cite{strogatz2015,murray2002}. La bifurcación de Hopf es la vía canónica por la cual un estado estacionario pierde su estabilidad y engendra espontáneamente una oscilación periódica autosostenida (un ciclo límite).

\subsection{El Teorema de Hopf en el plano}

Consideremos un sistema bidimensional parametrizado por $\mu \in \mathbb{R}$:
\begin{equation}
    \dot{\vec{x}} = \vec{f}(\vec{x}; \mu), \quad \vec{x} \in \mathbb{R}^2.
\end{equation}
Supongamos que el sistema posee un punto fijo de equilibrio $\vec{x}^*(\mu)$ para todos los valores de $\mu$ cercanos a un valor crítico $\mu_c = 0$, de modo que $\vec{f}(\vec{x}^*(\mu); \mu) = \vec{0}$.

\begin{theorem}[Teorema de la bifurcación de Hopf (Andronov--Hopf)]
\label{thm:hopf_teorema}
Sea $J(\mu) = D\vec{f}(\vec{x}^*(\mu); \mu)$ la matriz jacobiana evaluada en el equilibrio. Supongamos que:
\begin{enumerate}
    \item Los autovalores de $J(\mu)$ son un par de complejos conjugados:
    \begin{equation}
        \lambda_{1,2}(\mu) = \alpha(\mu) \pm i\omega(\mu).
    \end{equation}
    \item En el valor crítico $\mu = 0$, los autovalores cruzan el eje imaginario:
    \begin{equation}
        \alpha(0) = 0, \qquad \omega_0 \equiv \omega(0) > 0.
    \end{equation}
    \item \textbf{Condición de transversalidad:} La parte real de los autovalores cruza el eje imaginario con velocidad estrictamente no nula:
    \begin{equation}
        \left. \frac{d\alpha}{d\mu} \right|_{\mu = 0} \neq 0.
    \end{equation}
\end{enumerate}
Entonces, al variar $\mu$ a través de $\mu = 0$, el punto fijo $\vec{x}^*(\mu)$ cambia de estabilidad (de foco estable a inestable o viceversa) y \textbf{nace o muere una familia de órbitas periódicas cerradas (ciclos límite)} en una vecindad del equilibrio, con frecuencia angular inicial $\omega_0$ y periodo $T_0 \approx 2\pi/\omega_0$.
\end{theorem}

Dependiendo del signo de los términos no lineales de orden superior en el punto crítico, la bifurcación de Hopf se clasifica en dos tipos fundamentales: \textbf{supercrítica} y \textbf{subcrítica}.

\subsection{Bifurcación de Hopf supercrítica}

En la bifurcación supercrítica, la transición hacia la oscilación es continua, suave y reversible (\emph{soft onset}).

\paragraph{Forma normal en coordenadas polares.}
La forma normal canónica en coordenadas polares $(r, \theta)$ se escribe:
\begin{align}
    \dot{r} &= \mu r - r^3, \label{eq:hopf_super_r} \\
    \dot{\theta} &= \omega + b r^2, \label{eq:hopf_super_theta}
\end{align}
donde $\mu$ es el parámetro de bifurcación, $\omega > 0$ es la frecuencia básica y $b \in \mathbb{R}$ modula el cambio de frecuencia con la amplitud. Observemos que la ecuación radial $\dot{r} = \mu r - r^3$ tiene formalmente la estructura de una bifurcación pitchfork supercrítica en el radio $r \ge 0$:

\begin{itemize}
    \item \textbf{Régimen $\mu < 0$:} El único equilibrio radial es $r^* = 0$ (el origen). Dado que $\dot{r} = \mu r - r^3 < 0$ para todo $r > 0$, todas las trayectorias espiralan convergiendo exponencialmente al origen. El origen es un \textbf{foco estable atractor}.
    
    \item \textbf{Punto crítico $\mu = 0$:} La ecuación radial se reduce a $\dot{r} = -r^3$. El origen sigue siendo asintóticamente estable, pero la relajación es lenta y algebraica ($r(t) \sim t^{-1/2}$).
    
    \item \textbf{Régimen $\mu > 0$:} En el origen, $\left.\frac{d\dot{r}}{dr}\right|_{r=0} = \mu > 0$, por lo que el equilibrio $r^* = 0$ se vuelve un \textbf{foco inestable repulsor}. Simultáneamente, aparece un nuevo equilibrio radial no trivial dado por $\mu r - r^3 = 0 \implies r^*(\mu) = \sqrt{\mu}$.
\end{itemize}

Como $\dot{\theta} \approx \omega > 0$, la solución en $r = \sqrt{\mu}$ corresponde a una trayectoria periódica circular cerrada de radio constante: un \textbf{ciclo límite estable} (véase la \cref{fig:hopf_supercritica}).

\begin{figure}[htbp]
\centering
\begin{tikzpicture}[>=Stealth, scale=0.9]
    % Panel izquierdo: Espacio de fases para mu < 0
    \begin{scope}[shift={(-4.5, 0)}]
        \draw[->, thick, black!60] (-2,0) -- (2,0) node[right, font=\tiny] {$x$};
        \draw[->, thick, black!60] (0,-2) -- (0,2) node[above, font=\tiny] {$y$};
        \filldraw[black] (0,0) circle (2.2pt);
        % Espiral hacia adentro
        \draw[thick, black!85, domain=0:9.42, samples=120, variable=\t]
            plot ({ 1.6*exp(-0.35*\t)*cos(\t r) }, { 1.6*exp(-0.35*\t)*sin(\t r) });
        \node[font=\footnotesize, align=center] at (0,-2.4) {\textbf{(a)} $\mu < 0$\\Foco estable};
    \end{scope}

    % Panel central: Espacio de fases para mu > 0
    \begin{scope}[shift={(0, 0)}]
        \draw[->, thick, black!60] (-2,0) -- (2,0) node[right, font=\tiny] {$x$};
        \draw[->, thick, black!60] (0,-2) -- (0,2) node[above, font=\tiny] {$y$};
        \draw[black, fill=white, thick] (0,0) circle (2.2pt);
        % Ciclo límite estable
        \draw[very thick, black!95] (0,0) circle (1.25cm);
        % Espiral interna
        \draw[thick, black!65, domain=0:8.0, samples=100, variable=\t]
            plot ({ (0.2 + 0.12*\t)*cos(\t r) }, { (0.2 + 0.12*\t)*sin(\t r) });
        % Espiral externa
        \draw[thick, black!65, domain=0:6.28, samples=100, variable=\t]
            plot ({ (1.9 - 0.10*\t)*cos(\t r) }, { (1.9 - 0.10*\t)*sin(\t r) });
        \node[font=\footnotesize, align=center] at (0,-2.4) {\textbf{(b)} $\mu > 0$\\Ciclo límite ($r=\sqrt{\mu}$)};
    \end{scope}

    % Panel derecho: Diagrama de bifurcación
    \begin{scope}[shift={(4.5, 0)}]
        \draw[->, thick, black!60] (-2,0) -- (2.2,0) node[right, font=\tiny] {$\mu$};
        \draw[->, thick, black!60] (0,0) -- (0,2.2) node[above, font=\tiny] {$r$};
        % Origen: estable para mu < 0, inestable para mu > 0
        \draw[very thick, black!90] (-1.8,0) -- (0,0);
        \draw[very thick, black!90, dashed] (0,0) -- (1.8,0);
        % Rama de ciclo límite r = sqrt(mu)
        \draw[very thick, black!95, domain=0:1.8, samples=60] plot (\x, {1.2*sqrt(\x)});
        \filldraw[black] (0,0) circle (2.2pt);
        \node[font=\footnotesize, align=center] at (1.1, 1.7) {Ciclo límite\\$r^* = \sqrt{\mu}$};
        \node[font=\footnotesize, align=center] at (0,-2.4) {\textbf{(c)} Diagrama de\\bifurcación};
    \end{scope}
\end{tikzpicture}
\caption{Bifurcación de Hopf supercrítica. (a) Para $\mu < 0$, el origen es un foco estable. (b) Para $\mu > 0$, el origen se vuelve inestable y nace un ciclo límite atractor. (c) Diagrama de bifurcación mostrando el crecimiento continuo de la amplitud como $r \sim \sqrt{\mu}$.}
\label{fig:hopf_supercritica}
\end{figure}

\subsection{Bifurcación de Hopf subcrítica e histéresis}

En la bifurcación de Hopf subcrítica, el término cúbico no lineal tiene signo opuesto ($+r^3$) y actúa como desestabilizador. Para modelar un sistema físicamente confinado se incluye un término estabilizador de orden superior (quíntico, $-r^5$):
\begin{align}
    \dot{r} &= \mu r + r^3 - r^5, \label{eq:hopf_sub_r} \\
    \dot{\theta} &= \omega + b r^2. \label{eq:hopf_sub_theta}
\end{align}

\paragraph{Dinámica y catástrofe dinámica (\emph{Hard onset}).}
El comportamiento cualitativo revela una rica estructura de biestabilidad:
\begin{itemize}
    \item Para $\mu < 0$: El origen $r = 0$ es un foco estable, pero está rodeado por un \textbf{ciclo límite inestable} de radio $r_1^* \approx \sqrt{-\mu}$. Este ciclo inestable define la frontera de la cuenca de atracción del origen: si una perturbación saca al sistema fuera de este ciclo, la trayectoria es expulsada hacia un \textbf{ciclo límite estable de gran amplitud} ($r_2^* \approx 1$).
    \item En $\mu = 0$: El ciclo límite inestable colisiona con el origen y lo destruye.
    \item Para $\mu > 0$: El origen es inestable y no existen oscilaciones pequeñas; el sistema experimenta un salto brusco e instantáneo (\emph{hard onset}) hacia la oscilación de gran amplitud.
\end{itemize}

Si una vez encendida la oscilación se reduce nuevamente $\mu$ por debajo de cero ($\mu < 0$), el sistema no se apaga en $\mu = 0$, sino que permanece en la rama de oscilación grande hasta alcanzar el punto de silla-nodo de ciclos límite en $\mu_c < 0$. Este fenómeno es la \textbf{histéresis oscilatoria}.

\begin{figure}[htbp]
\centering
\begin{tikzpicture}[>=Stealth, scale=0.95]
    \draw[->, thick, black!60] (-2.5,0) -- (2.5,0) node[right, font=\small] {$\mu$};
    \draw[->, thick, black!60] (0,0) -- (0,2.6) node[above, font=\small] {$r$};

    % Eje r=0: Estable para mu < 0, inestable para mu > 0
    \draw[very thick, black!90] (-2.2,0) -- (0,0);
    \draw[very thick, black!90, dashed] (0,0) -- (2.2,0);
    \filldraw[black] (0,0) circle (2.2pt);

    % Ciclo límite inestable (línea discontinua) que nace subcríticamente hacia la izquierda
    \draw[thick, black!85, dashed, domain=-1.2:0, samples=60] plot (\x, {1.1*sqrt(-\x)});

    % Ciclo límite estable de gran amplitud (línea sólida superior)
    \draw[very thick, black!95, domain=-1.2:2.2, samples=80] plot (\x, {1.7 + 0.15*\x});

    % Pliegue en mu_c = -1.2
    \draw[thick, black!85, dashed] (-1.2, 1.2) -- (-1.2, 1.52);

    % Flechas de salto por histéresis
    \draw[->, thick, black!80] (0.05, 0.1) -- (0.05, 1.6);
    \node[right, font=\tiny] at (0.05, 0.85) {Salto brusco (\emph{hard onset})};

    \draw[->, thick, black!80] (-1.25, 1.4) -- (-1.25, 0.1);
    \node[left, font=\tiny] at (-1.25, 0.75) {Colapso};

    \node[font=\footnotesize, align=center] at (-1.5, 0.8) {Ciclo inestable\\(frontera cuenca)};
    \node[font=\footnotesize, align=center] at (1.2, 2.1) {Ciclo estable\\de gran amplitud};
    \node[font=\footnotesize, align=center] at (-1.4, -0.4) {Biestabilidad};
\end{tikzpicture}
\caption{Diagrama de bifurcación de Hopf subcrítica. La rama inestable (línea discontinua) retrocede hacia $\mu < 0$ y se pliega conectando con un ciclo límite estable de gran amplitud, generando histéresis y biestabilidad entre el reposo y la oscilación.}
\label{fig:hopf_subcritica}
\end{figure}

\subsection{Aplicaciones biológicas: excitabilidad neuronal y relojes bioquímicos}

\paragraph{Excitabilidad neuronal Tipo II (Hodgkin--Huxley).}
En el célebre modelo electrofisiológico del axón gigante de calamar formulado por Alan Hodgkin y Andrew Huxley (1952), el potencial de membrana $V(t)$ y las variables de compuerta iónica de sodio y potasio $(m, n, h)$ responden a una corriente externa inyectada $I_{\text{ext}}$.

Para corrientes bajas ($I_{\text{ext}} < I_c$), la neurona permanece en un estado de reposo hiperpolarizado estable. Al superar la corriente crítica ($I_{\text{ext}} > I_c$), el equilibrio de reposo pierde estabilidad a través de una \textbf{bifurcación de Hopf subcrítica}. El potencial salta bruscamente a un tren de potenciales de acción periódicos (\emph{spiking}) con una frecuencia de disparo no nula y finita desde el encendido. Esta respuesta se conoce en neurociencia como \textbf{excitabilidad neuronal de Tipo II}~\cite{edelstein2007}.

\paragraph{Oscilaciones glicolíticas y relojes genéticos.}
En la bioquímica celular, las oscilaciones periódicas en la concentración de intermediarios metabólicos como el NADH en levaduras y las oscilaciones sintéticas en circuitos genéticos (como el \emph{Repressilator}) emergen cuando la tasa de sustrato o la inhibición por retroalimentación negativa supera el umbral crítico establecido por el teorema de Hopf\aside{Para detectar numéricamente bifurcaciones de Hopf en sistemas no lineales de alta dimensión, se emplean algoritmos de continuación numérica como AUTO o XPPAUT, que rastrean los ceros de $\operatorname{Re}(\lambda)$ a lo largo de ramas de equilibrios.}.
